\documentclass[10pt,a4paper]{amsart}
\usepackage{amsmath,amssymb,amsthm,graphicx}
% For commutative diagrams we use Xy-pic
\usepackage[all]{xy}

% We use 2cell for 2-commutative diagrams.
\xyoption{2cell}
% Theorem environments.
%
\theoremstyle{plain}
\newtheorem{theorem}[subsection]{Theorem}
\newtheorem{proposition}[subsection]{Proposition}
\newtheorem{lemma}[subsection]{Lemma}

\theoremstyle{definition}
\newtheorem{definition}[subsection]{Definition}
\newtheorem{example}[subsection]{Example}
\newtheorem{exercise}[subsection]{Exercise}
\newtheorem{situation}[subsection]{Situation}

\theoremstyle{remark}
\newtheorem{remark}[subsection]{Remark}
\newtheorem{remarks}[subsection]{Remarks}

\numberwithin{equation}{subsection}

% Macros
%
\def\lim{\mathop{\mathrm{lim}}\nolimits}
\def\colim{\mathop{\mathrm{colim}}\nolimits}
\def\Spec{\mathop{\mathrm{Spec}}}
\def\Hom{\mathop{\mathrm{Hom}}\nolimits}
\def\Ext{\mathop{\mathrm{Ext}}\nolimits}
\def\SheafHom{\mathop{\mathcal{H}\!\mathit{om}}\nolimits}
\def\SheafExt{\mathop{\mathcal{E}\!\mathit{xt}}\nolimits}
\def\Sch{\mathit{Sch}}
\def\Mor{\mathop{\mathrm{Mor}}\nolimits}
\def\Ob{\mathop{\mathrm{Ob}}\nolimits}
\def\Sh{\mathop{\mathit{Sh}}\nolimits}
\def\NL{\mathop{N\!L}\nolimits}
\def\CH{\mathop{\mathrm{CH}}\nolimits}
\def\proetale{{pro\text{-}\acute{e}tale}}
\def\etale{{\acute{e}tale}}
\def\QCoh{\mathit{QCoh}}
\def\Ker{\mathop{\mathrm{Ker}}}
\def\Im{\mathop{\mathrm{Im}}}
\def\Coker{\mathop{\mathrm{Coker}}}
\def\Coim{\mathop{\mathrm{Coim}}}

% Boxtimes
%
\DeclareMathSymbol{\boxtimes}{\mathbin}{AMSa}{"02}

%
% Macros for moduli stacks/spaces
%
\def\QCohstack{\mathcal{QC}\!\mathit{oh}}
\def\Cohstack{\mathcal{C}\!\mathit{oh}}
\def\Spacesstack{\mathcal{S}\!\mathit{paces}}
\def\Quotfunctor{\mathrm{Quot}}
\def\Hilbfunctor{\mathrm{Hilb}}
\def\Curvesstack{\mathcal{C}\!\mathit{urves}}
\def\Polarizedstack{\mathcal{P}\!\mathit{olarized}}
\def\Complexesstack{\mathcal{C}\!\mathit{omplexes}}
% \Pic is the operator that assigns to X its picard group, usage \Pic(X)
% \Picardstack_{X/B} denotes the Picard stack of X over B
% \Picardfunctor_{X/B} denotes the Picard functor of X over B
\def\Pic{\mathop{\mathrm{Pic}}\nolimits}
\def\Picardstack{\mathcal{P}\!\mathit{ic}}
\def\Picardfunctor{\mathrm{Pic}}
\def\Deformationcategory{\mathcal{D}\!\mathit{ef}}
\usepackage{cleveref}
\crefname{originallemma}{original lemma}{original lemmas}
\Crefname{originallemma}{Original lemma}{Original lemmas}
\makeatletter
\newlabel{sites-cohomology-lemma-restrict-direct-image-open}{{0D6G}{source}{}{sourcepointer-1}{}}
\newlabel{sites-cohomology-lemma-restrict-direct-image-open@cref}{{[originallemma][0][]0D6G}{[1][0][]source}{}{}{}}
\makeatother
\crefname{originallemma}{original lemma}{original lemmas}
\Crefname{originallemma}{Original lemma}{Original lemmas}
\makeatletter
\newlabel{sites-lemma-localize-morphism-strong}{{03EF}{source}{}{sourcepointer-2}{}}
\newlabel{sites-lemma-localize-morphism-strong@cref}{{[originallemma][0][]03EF}{[1][0][]source}{}{}{}}
\makeatother
\title{Base change statements can be checked after an etale covering that factors through etale morphisms to the base}
\author{The Stacks Project Authors}
\date{stacks-project:etale-cohomology.tex:local-snapshot-20261011; GNU FDL 1.2 or later}
\begin{document}
\maketitle
\noindent\textit{Editorial note.} All delivered mathematical text is the author's own, copied from the pinned Stacks Project chapter file etale-cohomology.tex; no statement or proof step was rewritten, reordered or removed.\par\medskip
\noindent\textit{Editorial note.} The counted proof is the author's complete argument: the proof is written out in full for part (1) and then transported to parts (2) and (3) by replacing the locality input; the two displayed computations of the restrictions are given with the reason for each equality.\par\medskip
\noindent\textit{Editorial note.} The label lemma-base-change-local names the counted statement, and the proof environment immediately following it in the source is delivered verbatim as its proof.\par\medskip
\noindent\textit{Editorial note.} The counted statement is carried by the author's own lemma environment in the chapter source.\par\medskip
\noindent\textit{Editorial note.} The statement's two displayed cartesian diagrams are the author's and are delivered verbatim; the two cited inputs are kept as external pointers with their canonical Stacks tags.\par\medskip
\noindent\textit{Editorial note.} The preamble is the chapter's own native preamble excerpt already used for this source: the Xy-pic setup (native preamble lines 14-18) and the native theorem declarations with the mathematical macros (native preamble lines 154-217).\par\medskip
\noindent\textit{Allowed standard prerequisites (editorial).}\begin{itemize}
\item Locality of the pushforward and pullback along a localisation. Sites, the lemma computing the restriction of a pushforward and of a pullback along an etale morphism to a localisation. Supplies the two computations g\_*F|V\_i = g\_\{i,*\}F\_i and h\_*e\textasciicircum{}\{-1\}F|U\_i = h\_\{i,*\}e\_i\textasciicircum{}\{-1\}F\_i.
\item The same locality for the higher direct images. Cohomology on Sites, the lemma restricting higher direct images to an open. Supplies the analogous computations for the derived and abelian cases.
\end{itemize}\medskip
\ifdefined\hypertarget\hypertarget{sourcepointer-1}{}\fi
\noindent\textit{Source reference.} Original lemma 0D6G: sources/stacks-chow/tags.txt, source lines 16521--16521. Prerequisite pointer invoked at proof line(s) 16369; it names a standard fact of a sibling chapter that the counted argument uses, and every step of the counted proof built on it is retained verbatim in the delivered body.\par\medskip
\ifdefined\hypertarget\hypertarget{sourcepointer-2}{}\fi
\noindent\textit{Source reference.} Original lemma 03EF: sources/stacks-chow/tags.txt, source lines 4454--4454. Prerequisite pointer invoked at proof line(s) 16349; it names a standard fact of a sibling chapter that the counted argument uses, and every step of the counted proof built on it is retained verbatim in the delivered body.\par\medskip
\begin{lemma}
\label{lemma-base-change-local}
Consider a cartesian diagram of schemes
$$
\xymatrix{
X \ar[d]_f & Y \ar[l]^h \ar[d]^e \\
S & T \ar[l]_g
}
$$
Let $\{U_i \to X\}$ be an \'etale covering such that $U_i \to S$
factors as $U_i \to V_i \to S$ with $V_i \to S$ \'etale
and consider the cartesian diagrams
$$
\xymatrix{
U_i \ar[d]_{f_i} & U_i \times_X Y \ar[l]^{h_i} \ar[d]^{e_i} \\
V_i & V_i \times_S T \ar[l]_{g_i}
}
$$
Let $\mathcal{F}$ be a sheaf on $T_\etale$. Let $K$ in $D(T_\etale)$.
Set $K_i = K|_{V_i \times_S T}$ and
$\mathcal{F}_i = \mathcal{F}|_{V_i \times_S T}$.
\begin{enumerate}
\item If $f_i^{-1}g_{i, *}\mathcal{F}_i = h_{i, *}e_i^{-1}\mathcal{F}_i$
for all $i$, then $f^{-1}g_*\mathcal{F} = h_*e^{-1}\mathcal{F}$.
\item If $f_i^{-1}Rg_{i, *}K_i = Rh_{i, *}e_i^{-1}K_i$
for all $i$, then $f^{-1}Rg_*K = Rh_*e^{-1}K$.
\item If $\mathcal{F}$ is an abelian sheaf and
$f_i^{-1}R^qg_{i, *}\mathcal{F}_i = R^qh_{i, *}e_i^{-1}\mathcal{F}_i$
for all $i$, then
$f^{-1}R^qg_*\mathcal{F} = R^qh_*e^{-1}\mathcal{F}$.
\end{enumerate}
\end{lemma}

\begin{proof}
Proof of (1). First we observe that
$$
(f^{-1}g_*\mathcal{F})|_{U_i} =
f_i^{-1}(g_*\mathcal{F}|_{V_i}) =
f_i^{-1}g_{i, *}\mathcal{F}_i
$$
The first equality because $U_i \to X \to S$ is equal to
$U_i \to V_i \to S$ and the second equality because
$g_*\mathcal{F}|_{V_i} = g_{i, *}\mathcal{F}_i$ by
Sites, Lemma \ref{sites-lemma-localize-morphism-strong}.
Similarly we have
$$
(h_*e^{-1}\mathcal{F})|_{U_i} =
h_{i, *}(e^{-1}\mathcal{F}|_{U_i \times_X Y}) =
h_{i, *}e_i^{-1}\mathcal{F}_i
$$
Thus if the base change maps
$f_i^{-1}g_{i, *}\mathcal{F}_i \to h_{i, *}e_i^{-1}\mathcal{F}_i$
are isomorphisms for all $i$, then the base change map
$f^{-1}g_*\mathcal{F} \to h_*e^{-1}\mathcal{F}$
restricts to an isomorphism over $U_i$ for all $i$
and we conclude it is an isomorphism as $\{U_i \to X\}$
is an \'etale covering.

\medskip\noindent
For the other two statements we replace the appeal to
Sites, Lemma \ref{sites-lemma-localize-morphism-strong}
by an appeal to
Cohomology on Sites, Lemma
\ref{sites-cohomology-lemma-restrict-direct-image-open}.
\end{proof}

\end{document}
